% !PRINTED 701
\weberpage{701}
%<<CONT
ist. Ein Differential, dessen Untereck nur aus zwei einfachen Punkten besteht, heißt ein Differential dritter Gattung. Wir beweisen den Satz:

4. \gesperrt{Jedes Differential $dJ$ läßt sich linear und mit konstanten Koeffizienten aus Differentialen erster, zweiter und dritter Gattung zusammensetzen.}

Dieser Satz kann nach (4) auch so ausgesprochen werden, daß die Klasse $A$ eine Basis
\weq{6}{\mathfrak{A}_1, \mathfrak{A}_2', \ldots, \mathfrak{A}_{b+p-1}}
von der Art hat, daß die Quotienten
\weq{7}{dJ_r = \frac{\mathfrak{A}_r}{\mathfrak{B}}}
Differentiale von einer der drei Gattungen sind.

Um ihn durch vollständige Induktion zu beweisen, nehmen wir an, es sei für irgend ein $\mathfrak{B}$ eine solche Basis gefunden, und suchen eine ebensolche für das Untereck $\mathfrak{P}\mathfrak{B}$, wenn $\mathfrak{P}$ ein beliebiger Punkt ist. Eine solche Basis können wir in der Form annehmen:
\weq{8}{\mathfrak{P}\mathfrak{A}_1, \mathfrak{P}\mathfrak{A}_2, \ldots, \mathfrak{P}\mathfrak{A}_{b+p-1}, \mathfrak{A}_{b+p},}
worin $\mathfrak{A}_{b+p}$ ein mit $\mathfrak{W}\mathfrak{B}\mathfrak{P}$ äquivalentes Polygon sein muß, das den Punkt $\mathfrak{P}$ nicht enthält.

Wenn der Punkt $\mathfrak{P}$ in $\mathfrak{B}$ aufgeht, setzen wir:
\[
\mathfrak{B} = \mathfrak{M}\mathfrak{P}^m,
\]
so daß $\mathfrak{P}$ nicht mehr in $\mathfrak{M}$ aufgeht. Das Polygon $\mathfrak{W}\mathfrak{P}^{m+1}$ gehört dann in eine eigentliche Klasse, in der also ein durch $\mathfrak{P}$ nicht teilbares Polygon $\mathfrak{N}$ existiert.

Wir setzen dann:
\[
\mathfrak{A}_{b+p} = \mathfrak{M}\mathfrak{N},
\]
und erhalten daraus ein Differential zweiter Gattung:
\weq{9}{dJ_{b+p} = \frac{\mathfrak{A}_{b+p}}{\mathfrak{B}\mathfrak{P}} = \frac{\mathfrak{N}}{\mathfrak{P}^{m+1}}.}

Geht $\mathfrak{P}$ in $\mathfrak{B}$ nicht auf, so nehme man einen in $\mathfrak{B}$ aufgehenden Punkt $\mathfrak{P}_0$. Dann gehört $\mathfrak{W}\mathfrak{P}_0\mathfrak{P}$ in eine eigentliche Klasse und enthält ein durch $\mathfrak{P}$ nicht teilbares Polygon $\mathfrak{N}$. Ist dann
\[
\mathfrak{B} = \mathfrak{P}_0\mathfrak{M}, \quad \mathfrak{A}_{b+p} = \mathfrak{M}\mathfrak{N},
\]
so ist
\weq{10}{dJ_{b+p} = \frac{\mathfrak{A}_{b+p}}{\mathfrak{B}\mathfrak{P}} = \frac{\mathfrak{N}}{\mathfrak{P}_0\mathfrak{P}}}
%CONT>>
